\section{Applicability of the analytic results}
\label{sec:analytic-family}

The reconstruction theorem applies on the known period ball
$B_{\rm ev}$. To relate its output to the analytic results of
\cite{W_Hodge,W_HYM}, we verify their hypotheses for the
selected central member and its evaluated smoothing family.
Appendix~\ref{app:analytic-inputs} states those results in full,
including the central metrics, harmonic masses, bundle
parameters and spatial norms used below.

Fix a finite graph regularity order $m\ge4$. We shall obtain
a full positive-radius ball $B_{\rm an}\subset B_{\rm ev}$,
together with positive bundle-parameter and metric-basin radii.
These radii may depend on the selected central coefficients
and on $m$, and need not be computable. For
$p\in B_{\rm an}^{\circ}$, write $g_p$ for the Ricci-flat
K\"ahler metric in $c_1(\OO(2,2)|_{X_p})$.

\begin{proposition}\label{prop:analytic-reapplication}
The family of Theorem~\ref{thm:evaluated-periods} admits
$B_{\rm an}$ and the analytic objects specified in
Appendix~\ref{app:analytic-inputs}, satisfying
\eqref{eq:actual-metric-input} and
Recollections~\ref{thm:hodge-input}--\ref{thm:moduli-input}.
The observation procedure of
Theorem~\ref{thm:integral-reconstruction} remains defined on
all of $B_{\rm ev}^{\circ}$ and in particular applies to
each of these analytic target fibres.
\end{proposition}

\begin{proof}
First, the selected central transition is linked to the original
general one within the same fixed toric construction. Let
$\mathcal H^{\rm geom}$ require only a smooth $Y_h$ and four
simple zeros on each singular line. It is a nonempty
Zariski-open subset of $\mathcal H$, and hence path connected.
It contains both the original general choice and $h_*$.
The universal resolved divisor is smooth and proper there.
Transport along a connecting path preserves the four ambient
divisor classes, the exact polarization and all intersections.
The four roots on each fixed line can be followed along that
path, so the two exceptional types remain distinct and unchanged,
up to permutations within each group. This identifies the
topological inputs, not the central metrics or analytic constants.

The selected $X_0$ is normal Gorenstein Calabi--Yau with a
projective small resolution, the required structure-sheaf
vanishings, and a polarized smoothing. A line $p=tv$ with
every $v_i\ne0$ supplies the smoothing disc needed for the
fixed-central conical metric theorem used in
\cite[Section~7]{W_Hodge}. Choose that metric $g_0$ anew.
The evaluated Morse coordinates need not be its conical
coordinates. The coordinate reconciliation in the same proof
uses a central ambient map with $Q\circ\Phi=aQ$; replacing
$\Phi$ by $a^{-1/2}\Phi$ makes it preserve $Q$ and its nearby
quadric fibres. The residual coordinate correction has positive
conic order. These analytic charts coexist with the evaluated
charts; the observation procedure does not evaluate the
conical coordinate change.

The metric construction needs this fixed central conical metric,
a smooth proper exterior, the relative polarization and form,
and holomorphic quadric charts with bounded units. All are
available here. Choose anew its central inverses, normalizations
and finite low-growth modes. The independent-cap inverse and
nonlinear correction in \cite[Sections~7--9]{W_Hodge} then
give \eqref{eq:actual-metric-input} for the actual polarized
Ricci-flat metrics. Their uniqueness identifies them with $g_p$.
The Hodge trials, capacity lower bounds and mass compactness
use this comparison and the same exterior restriction map.
The primitive K\"ahler mass identity and
\eqref{eq:central-intersection}, with the full surviving
space removed, give exactly $M_H$ in \eqref{eq:hodge-model}. The complete
degree-one and degree-two argument in \cite[Sections~10--12]{W_Hodge}
therefore gives Recollection~\ref{thm:hodge-input} on these fibres.
No uniform estimate over varying $h$ follows from this reasoning.

The bundle construction requires an additional ambient
fourfold condition. At $c_*$ the first two equations form
a regular sequence with normal integral zero set $V_0$,
and $h_*$ is a nonzerodivisor. These properties persist on
an open coefficient neighbourhood: the complete-intersection
family is flat, normal geometric integrality is open, and
the regular-sequence conditions are open. Shrink the analytic
ball so that $\sigma\pi_R$ lies in this neighbourhood.
Adjunction then gives $K_{V_s}=\OO(-2,-2)$ and makes $X_s$
a relative anticanonical Cartier divisor. Smoothness of the
triple alone was not used to assert these pair properties.

The fixed-bidegree Koszul complexes in
\cite[Section~2]{W_HYM} now apply on this relative fourfold.
Their cohomology gives the stated extension dimensions,
vanishing ambient obstruction space and restriction
isomorphisms. It constructs the actual relative bundle and
the Chern data \eqref{eq:nonflat-characteristic-data}, with
the two Koszul frames chosen at $h_*$. The required equality
of the line Hilbert polynomials follows from these same
bidegrees and $H=\OO(2,2)$, with no symmetry assumption on $h_*$.

The HE line background and complementary estimates in
\cite[Sections~3--5]{W_HYM} use the actual metrics just
constructed, fixed-exterior estimates, the ambient metric
comparison, $H^0(L_p^{\pm2})=0$, and Hartogs extension at
the nodes. The nonzero exterior restriction of $\pm2D_L$
follows from the resolution Thom sequence and $D_L^2H=-16$;
simple connectivity alone is not used to exclude off-diagonal
parallel sections. The supported extension representatives
come from total node trivializations and the fixed Koszul
classes. Reconstruct their central elliptic solutions and
positive masses $c_{\pm,0}$, and their harmonic continuations
$c_\pm(p)$. These constants are not imported numerically
from a different central member.

The local spin bound, critical smallness of the bounded
line connection and central extension computation give the
same all-degree split Dolbeault gaps. Together with the
section inverse they supply the projected HE correction
on a common product of base and matrix balls. The subsequent
real equivariant normalization, separate orbit balancing,
semistability and polystable coverage use precisely these
uniform estimates and positive masses. This proves that all
the hypotheses of Recollection~\ref{thm:bundle-input} hold here.

For flow, the required spatial heat bounds, connection
expansion and contracted divergence bounds come from this
same corrected family. The entry and confinement proof,
invariant graph and full-data phase in
\cite[Sections~9--11]{W_HYM} can consequently be constructed
on it, with the chosen finite $m$. Their smallness conditions
use the complementary gap, not the weakest positive finite
rate. The full $H^1$ gauge comparison, arbitrary-unitary
rigidity and actual harmonic class identification in
\cite[Sections~12--13]{W_HYM} apply to the same connections.
The finite balanced quotient and its confined paths then
give Recollection~\ref{thm:moduli-input}. Intersecting these
finitely many smallness conditions leaves full positive
base and bundle balls.

Finally $B_{\rm ev}$ retains its known radius and its
observation construction. The smaller $B_{\rm an}$ need
not contain the algebraic reference point or the chosen
meridian centres. Those remain in $B_{\rm ev}^{\circ}$:
their transports concern smooth complete intersections
and the same residue form. At a target $p\in B_{\rm an}^\circ$
they produce cycles and periods on that very $X_p$.
No PDE estimate is needed along the entire observation
path, and no computation of the radius of $B_{\rm an}$
is required to run the procedure on $B_{\rm ev}^\circ$.
\end{proof}
